% from ega0-1-fr.tex
\oldpage[0\textsubscript{I}]{43}
le vérifie sans peine~; on prend $v_\theta^\flat={v'}^\flat$. De même, de
$u:\mathcal{G}\to\Psi_*(\mathcal{F})$, qui est un homomorphisme de
$\mathcal{B}$-Modules, on déduit (3.7.1) un homomorphisme
$u^\sharp:\psi^*(\mathcal{G})\to\mathcal{F}$ de
$\psi^*(\mathcal{B})$-Modules, d'où par tensorisation avec $\mathcal{A}$ un
homomorphisme de $\mathcal{A}$-Modules
$\Psi^*(\mathcal{G})\to\mathcal{F}$, que nous désignerons par
$u_\theta^\sharp$. Il est immédiat de vérifier que
$(u_\theta^\sharp)_\theta^\flat=u$ et
$(v_\theta^\flat)_\theta^\sharp=v$, ainsi que le caractère fonctoriel en
$\mathcal{F}$ de l'isomorphisme $v\mapsto v_\theta^\flat$. Le caractère
fonctoriel en $\mathcal{G}$ de $u\mapsto u_\theta^\sharp$ s'en déduit alors
formellement comme dans (3.5.4) (raisonnement qui démontrerait aussi le
caractère fonctoriel de $\Psi^*$ établi dans (4.3.1) directement).

Si on prend pour $v$ l'homomorphisme identique de $\Psi^*(\mathcal{G})$,
$v_\theta^\flat$ est un homomorphisme
\[
  \rho_{\mathcal{G}}:\mathcal{G}\longrightarrow
  \Psi_*(\Psi^*(\mathcal{G}))~;
  \tag{4.4.3.2}\label{0.4.4.3.2-fr}
\]
si on prend pour $u$ l'homomorphisme identique de $\Psi_*(\mathcal{F})$,
$u_\theta^\sharp$ est un homomorphisme
\[
  \sigma_{\mathcal{F}}:\Psi^*(\Psi_*(\mathcal{F}))
  \longrightarrow\mathcal{F}~;
  \tag{4.4.3.3}\label{0.4.4.3.3-fr}
\]
ces homomorphismes seront dits \emph{canoniques}. Ils ne sont en général ni
injectifs ni surjectifs. On a des factorisations canoniques analogues à
(3.5.3.3) et (3.5.4.4).

On notera que si $s$ est une section de $\mathcal{G}$ au-dessus d'un ouvert
$V$ de $Y$, $\rho_{\mathcal{G}}(s)$ est la section $s'\otimes1$ de
$\Psi^*(\mathcal{G})$ au-dessus de $\psi^{-1}(V)$, $s'$ étant telle que
$s'_x=s_{\psi(x)}$ pour tout $x\in\psi^{-1}(V)$. Notons aussi que si
$u:\mathcal{G}\to\psi_*(\mathcal{F})$ est un homomorphisme, il définit pour
tout $x\in X$ un homomorphisme
$u_x:\mathcal{G}_{\psi(x)}\to\mathcal{F}_x$ sur les fibres, obtenu en
composant
$(u^\sharp)_x:(\Psi^*(\mathcal{G}))_x\to\mathcal{F}_x$ et l'homomorphisme
canonique $s_x\mapsto s_x\otimes1$ de $\mathcal{G}_{\psi(x)}$ dans
$(\Psi^*(\mathcal{G}))_x
=\mathcal{G}_{\psi(x)}\otimes_{\mathcal{B}_{\psi(x)}}\mathcal{A}_x$.
L'homomorphisme $u_x$ s'obtient aussi par passage à la limite inductive à
partir des homomorphismes
\[
  \Gamma(V,\mathcal{G})\xrightarrow{u}
  \Gamma(\psi^{-1}(V),\mathcal{F})\longrightarrow\mathcal{F}_x,
\]
où $V$ parcourt les voisinages de $\psi(x)$.
\end{env}

\begin{env}[4.4.4]
\label{0.4.4.4-fr}
Soient $\mathcal{F}_1$, $\mathcal{F}_2$ des $\mathcal{A}$-Modules,
$\mathcal{G}_1$, $\mathcal{G}_2$ des $\mathcal{B}$-Modules, $u_i$ ($i=1,2$)
un homomorphisme de $\mathcal{G}_i$ dans $\mathcal{F}_i$. Nous désignerons par
$u_1\otimes u_2$ l'homomorphisme
\[
  u:\mathcal{G}_1\otimes_{\mathcal{B}}\mathcal{G}_2
  \longrightarrow
  \mathcal{F}_1\otimes_{\mathcal{A}}\mathcal{F}_2
\]
tel que $u^\sharp=(u_1)^\sharp\otimes(u_2)^\sharp$ (compte tenu de
(4.3.3.1))~; on vérifie que $u$ est aussi le composé
\[
  \mathcal{G}_1\otimes_{\mathcal{B}}\mathcal{G}_2
  \longrightarrow
  \Psi_*(\mathcal{F}_1)\otimes_{\mathcal{B}}\Psi_*(\mathcal{F}_2)
  \longrightarrow
  \Psi_*(\mathcal{F}_1\otimes_{\mathcal{A}}\mathcal{F}_2),
\]
où la première flèche est le produit tensoriel ordinaire
$u_1\otimes_{\mathcal{B}}u_2$ et la seconde l'homomorphisme canonique
(4.2.2.1).
\end{env}

\begin{env}[4.4.5]
\label{0.4.4.5-fr}
Soient $(\mathcal{G}_\lambda)_{\lambda\in L}$ un système inductif de
$\mathcal{B}$-Modules, et, pour tout $\lambda\in L$, soit $u_\lambda$ un
homomorphisme $\mathcal{G}_\lambda\to\Psi_*(\mathcal{F})$, formant un système
inductif~; posons
$\mathcal{G}=\varinjlim\mathcal{G}_\lambda$ et
$u=\varinjlim u_\lambda$~; alors les $(u_\lambda)^\sharp$ forment un système
inductif d'homomorphismes $\Psi^*(\mathcal{G}_\lambda)\to\mathcal{F}$, et la
limite inductive de ce système n'est autre que $u^\sharp$.
\end{env}

\begin{env}[4.4.6]
\label{0.4.4.6-fr}
Soient $\mathcal{M}$, $\mathcal{N}$ deux $\mathcal{B}$-Modules, $V$ un ouvert
de $Y$, $U=\psi^{-1}(V)$~; l'application $v\mapsto\Psi^*(v)$ est un
homomorphisme
\[
  \operatorname{Hom}_{\mathcal{B}|V}(\mathcal{M}|V,\mathcal{N}|V)
  \longrightarrow
  \operatorname{Hom}_{\mathcal{A}|U}
  (\Psi^*(\mathcal{M})|U,\Psi^*(\mathcal{N})|U)
\]
\begin{sloppypar}
pour les structures de $\Gamma(V,\mathcal{B})$-module
($\operatorname{Hom}_{\mathcal{A}|U}
(\Psi^*(\mathcal{M})|U,\allowbreak\Psi^*(\mathcal{N})|U)$ est normalement muni d'une
structure de $\Gamma(U,\psi^*(\mathcal{B}))$-module, et grâce à
l'homomorphisme
